
    %%%%%%%%%%%%%%%%%%%%%%%%%%%%%%%%%%%%%%%%%%%%%%%%%%%%%%%%%%%%%%%%%%%%%%%%%%%%%%%%
    %%% Classe do documento
    \documentclass[12pt, addpoints]{exam}
    \UseRawInputEncoding

    %%%%%%%%%%%%%%%%%%%%%%%%%%%%%%%%%%%%%%%%%%%%%%%%%%%%%%%%%%%%%%%%%%%%%%%%%%%%%%%%
    % preambulo.tex
    %
    % Arquivo de configuração de packages para uso o LaTeX, conforme minhas
    % preferências e modelos pessoais.
    %
    % Para maiores informações, visite:
    %    https://github.com/abrantesasf/latex
    %
    % NÃO ALTERE SE NÃO SOUBER O QUE ESTÁ FAZENDO!
    %%%%%%%%%%%%%%%%%%%%%%%%%%%%%%%%%%%%%%%%%%%%%%%%%%%%%%%%%%%%%%%%%%%%%%%%%%%%%%%%


    %%%%%%%%%%%%%%%%%%%%%%%%%%%%%%%%%%%%%%%%%%%%%%%%%%%%%%%%%%%%%%%%%%%%%%%%%%%%%%%%
    %%% Estruturas de controle:
    \input{static/utils/controles}

    %%%%%%%%%%%%%%%%%%%%%%%%%%%%%%%%%%%%%%%%%%%%%%%%%%%%%%%%%%%%%%%%%%%%%%%%%%%%%%%%
    %%% Compilação condicional em PDF:
    \input{static/utils/pdf}

    %%%%%%%%%%%%%%%%%%%%%%%%%%%%%%%%%%%%%%%%%%%%%%%%%%%%%%%%%%%%%%%%%%%%%%%%%%%%%%%%
    %%% Configurações de layout da página:
    \input{static/utils/layout}

    %%%%%%%%%%%%%%%%%%%%%%%%%%%%%%%%%%%%%%%%%%%%%%%%%%%%%%%%%%%%%%%%%%%%%%%%%%%%%%%%
    %%% Configurações para autores:
    \input{static/utils/autores}

    %%%%%%%%%%%%%%%%%%%%%%%%%%%%%%%%%%%%%%%%%%%%%%%%%%%%%%%%%%%%%%%%%%%%%%%%%%%%%%%%
    %%% Cabeçalho e rodapé:
    %%% Atenção! É MUITO DIFÍCIL ajustar apropriadamente os running headers
    %%% e running footers da primeira vez. Você pode confiar no LaTeX e deixar que
    %%% ele faça o ajuste automaticamente ou pode quebrar a cabeça e
    %%% tentar melhorar algo que o LaTeX já faz muito bem, mesmo que não fique
    %%% exatamente do jeito que você quer. O que você prefere? Se quiser ajustar
    %%% manualmente, descomente a linha abaixo e faça as configurações no arquivo
    %%% "utils/headings.tex".
    %\input{static/utils/headings}

    %%%%%%%%%%%%%%%%%%%%%%%%%%%%%%%%%%%%%%%%%%%%%%%%%%%%%%%%%%%%%%%%%%%%%%%%%%%%%%%%
    %%% Configuração para as fontes do documento:
    %%% Se você quiser usar fontes próprias ou do sistema, precisará ajustar as
    %%% configurações no arquivo "utils/fontes.tex".
    %%% ATENÇÃO: por padrão eu utilizo XeLaTeX com fontes proprietárias projetadas
    %%% por Matthew Butterick, adquiridas em https://mbtype.com, que não podem ser
    %%% distribuídas devido à licença de utilização. Se você usar XeLaTeX, você
    %%% DEVERÁ OBRIGATORIAMENTE ajustar as fontes no arquivo "utils/fontes.tex".
    \input{static/utils/fontes}

    %%%%%%%%%%%%%%%%%%%%%%%%%%%%%%%%%%%%%%%%%%%%%%%%%%%%%%%%%%%%%%%%%%%%%%%%%%%%%%%%
    %%% Sumário:
    \input{static/utils/toc}

    %%%%%%%%%%%%%%%%%%%%%%%%%%%%%%%%%%%%%%%%%%%%%%%%%%%%%%%%%%%%%%%%%%%%%%%%%%%%%%%%
    %%% Matemática:
    \input{static/utils/matematica}

    %%%%%%%%%%%%%%%%%%%%%%%%%%%%%%%%%%%%%%%%%%%%%%%%%%%%%%%%%%%%%%%%%%%%%%%%%%%%%%%%
    %%% Notas de rodapé:
    \input{static/utils/footnotes}

    %%%%%%%%%%%%%%%%%%%%%%%%%%%%%%%%%%%%%%%%%%%%%%%%%%%%%%%%%%%%%%%%%%%%%%%%%%%%%%%%
    %%% Referências cruzadas, links, citações:
    \input{static/utils/crossrefs}

    %%%%%%%%%%%%%%%%%%%%%%%%%%%%%%%%%%%%%%%%%%%%%%%%%%%%%%%%%%%%%%%%%%%%%%%%%%%%%%%%
    %%% Configurações para os hiperlinks em PDF:
    \input{static/utils/pdfconfig}

    %%%%%%%%%%%%%%%%%%%%%%%%%%%%%%%%%%%%%%%%%%%%%%%%%%%%%%%%%%%%%%%%%%%%%%%%%%%%%%%%
    %%% Glossário e índice remissivo:
    \input{static/utils/glossindex}

    %%%%%%%%%%%%%%%%%%%%%%%%%%%%%%%%%%%%%%%%%%%%%%%%%%%%%%%%%%%%%%%%%%%%%%%%%%%%%%%%
    %%% Referências bibliográficas:
    \input{static/utils/bibliografia}

    %%%%%%%%%%%%%%%%%%%%%%%%%%%%%%%%%%%%%%%%%%%%%%%%%%%%%%%%%%%%%%%%%%%%%%%%%%%%%%%%
    %%% Cores:
    \input{static/utils/cores}

    %%%%%%%%%%%%%%%%%%%%%%%%%%%%%%%%%%%%%%%%%%%%%%%%%%%%%%%%%%%%%%%%%%%%%%%%%%%%%%%%
    %%% Computação:
    \input{static/utils/computacao}

    %%%%%%%%%%%%%%%%%%%%%%%%%%%%%%%%%%%%%%%%%%%%%%%%%%%%%%%%%%%%%%%%%%%%%%%%%%%%%%%%
    %%% Gráficos:
    \input{static/utils/graficos}

    %%%%%%%%%%%%%%%%%%%%%%%%%%%%%%%%%%%%%%%%%%%%%%%%%%%%%%%%%%%%%%%%%%%%%%%%%%%%%%%%
    %%% Ícones:
    \input{static/utils/icones}

    %%%%%%%%%%%%%%%%%%%%%%%%%%%%%%%%%%%%%%%%%%%%%%%%%%%%%%%%%%%%%%%%%%%%%%%%%%%%%%%%
    %%% Blocos:
    \input{static/utils/blocos}

    %%%%%%%%%%%%%%%%%%%%%%%%%%%%%%%%%%%%%%%%%%%%%%%%%%%%%%%%%%%%%%%%%%%%%%%%%%%%%%%%
    %%% Boxes coloridos:
    \input{static/utils/colorbox}

    %%%%%%%%%%%%%%%%%%%%%%%%%%%%%%%%%%%%%%%%%%%%%%%%%%%%%%%%%%%%%%%%%%%%%%%%%%%%%%%%
    %%% Tabelas:
    \input{static/utils/tabelas}

    %%%%%%%%%%%%%%%%%%%%%%%%%%%%%%%%%%%%%%%%%%%%%%%%%%%%%%%%%%%%%%%%%%%%%%%%%%%%%%%%
    %%% Ambientes de listas:
    \input{static/utils/listas}

    %%%%%%%%%%%%%%%%%%%%%%%%%%%%%%%%%%%%%%%%%%%%%%%%%%%%%%%%%%%%%%%%%%%%%%%%%%%%%%%%
    %%% Ambientes floats e similares:
    \input{static/utils/floats}

    %%%%%%%%%%%%%%%%%%%%%%%%%%%%%%%%%%%%%%%%%%%%%%%%%%%%%%%%%%%%%%%%%%%%%%%%%%%%%%%%
    %%% Ajustes para a classe EXAM:
    \input{static/utils/exam}

    %%%%%%%%%%%%%%%%%%%%%%%%%%%%%%%%%%%%%%%%%%%%%%%%%%%%%%%%%%%%%%%%%%%%%%%%%%%%%%%%
    %%% Ajustes para textos:
    \input{static/utils/texto}

    %%%%%%%%%%%%%%%%%%%%%%%%%%%%%%%%%%%%%%%%%%%%%%%%%%%%%%%%%%%%%%%%%%%%%%%%%%%%%%%%
    %%% Ambientes, aliases, comandos e customizações em geral para serem
    %%% utilizadas nos documentos. Defina aqui qualquer customização específica
    %%% para seus documentos!
    \input{static/utils/customizacoes}

    %%%%%%%%%%%%%%%%%%%%%%%%%%%%%%%%%%%%%%%%%%%%%%%%%%%%%%%%%%%%%%%%%%%%%%%%%%%%%%%%
    %%% Ajuste do layout, espaçamento de linhas e etc.:
    \geometry{a4paper, portrait, left=2.0cm, right=2.0cm, top=2.0cm, bottom=2.0cm}
    \singlespacing

    %%%%%%%%%%%%%%%%%%%%%%%%%%%%%%%%%%%%%%%%%%%%%%%%%%%%%%%%%%%%%%%%%%%%%%%%%%%%%%%%
    %%% Configurações de cabeçalho e rodapé (não use fancyhdr pois ocorre conflito):
    \pagestyle{headandfoot}
    \runningheadrule
    \firstpageheader{}{}{}
    \runningheader{Sem disciplina}{}{Avaliação}
    \firstpagefooter{}{}{Página \thepage\ de \numpages}
    \runningfooter{}{}{Página \thepage\ de \numpages}
    \runningfootrule

    %%%%%%%%%%%%%%%%%%%%%%%%%%%%%%%%%%%%%%%%%%%%%%%%%%%%%%%%%%%%%%%%%%%%%%%%%%%%%%%%
    %%% Configurações para as propriedades do PDF:
    \hypersetup{
    hidelinks,           % Comente para web, descomente para impressão
    colorlinks=true,      % True para web, False para impressão
    pdftitle=TaRodqando: Avaliação,
    pdfauthor=Matheus Endlich Silveira,
    pdfsubject=Sem disciplina,
    pdfkeywords=['Sem tag'],
    pdfinfo={
        CreationDate={}, % Ex.: D:AAAAMMDDHH24MISS
        ModDate={}       % Ex.: D:AAAAMMDDHH24MISS
    }
    }
    \input{static/utils/pdfconfig}

    %%%%%%%%%%%%%%%%%%%%%%%%%%%%%%%%%%%%%%%%%%%%%%%%%%%%%%%%%%%%%%%%%%%%%%%%%%%%%%%%
    %%% Começa o documento
    \begin{document}

    %%%%%%%%%%%%%%%%%%%%%%%%%%%%%%%%%%%%%%%%%%%%%%%%%%%%%%%%%%%%%%%%%%%%%%%%%%%%%%%%
    %%% Folha de instruções padronizadas da UVV para a prova
    \pdfbookmark[1]{Instruções}{instrucoes}
    \begin{figure}[H]
        \begin{center}
            \includegraphics[scale=0.3]{{static/imgs/logo-uvv.png}}
        \end{center}	
    \end{figure}

    \vspace{-1cm}

    \begin{table}[H]
        \centering
        \begin{tabular}{|llll|}
            \hline
            \multicolumn{2}{|l|}{\textbf{Disciplina:} Sem disciplina} & 
            \multicolumn{1}{l|}{\multirow{3}{*}{\textbf{\makecell{Nota:\\ \,\\ \,}}}} & 
            \multirow{3}{*}{\textbf{\makecell{Coordenador:\\ \,\\ \,}}} \\ 
            \cline{1-2}
            \multicolumn{2}{|l|}{\textbf{Professor:} Matheus Endlich Silveira \hspace{4.72cm} } & 
            \multicolumn{1}{l|}{} & \\ 
            \cline{1-2}
            \multicolumn{2}{|l|}{\textbf{Aluno:}} & 
            \multicolumn{1}{l|}{} & \\ 
            \hline
            \multicolumn{1}{|l|}{\textbf{Turma:} \hspace{6.3cm} } & 
            \multicolumn{1}{l|}{\textbf{Semestre:}} & 
            \multicolumn{2}{l|}{\textbf{Valor:} 10 pontos} \\ 
            \hline
            \multicolumn{1}{|l|}{\textbf{Data:}} & 
            \multicolumn{3}{l|}{\textbf{Avaliação:}} \\ 
            \hline
        \end{tabular}
    \end{table}

    \begin{center}
        \fbox{\setlength\fboxsep{0.3cm} \fbox{\parbox{15.4cm}{
            \begin{center}
                \textbf{INSTRUÇÕES PARA A REALIZAÇÃO DESTA AVALIAÇÃO:}
            \end{center}
            \begin{itemize}[leftmargin=*]
                \item Esta é a primeira avaliação da disciplina Sem disciplina, 
                    e engloba o conteúdo referente Sem tag.
                \item Esta avaliação contém 50 questões objetivas, \textbf{todas obrigatórias}.
                \item \textbf{Leia atentamente cada questão!} As questões objetivas terão uma,
                    e apenas uma única, resposta a ser marcada.
                \item \textbf{Desligue o celular} antes de começar. A avaliação é
                    \textbf{individual} e \textbf{sem consulta}.
                \item As questões podem ser respondidas, na prova, com lápis ou caneta. Ao final
                    da prova \textbf{{VOCÊ DEVERÁ PREENCHER O GABARITO COM CANETA
                    DE COR PRETA}} \textbf{\underline{OU AZUL}} para que seja feita a correção
                    automatizada através da leitura do padrão de respostas marcado no
                    gabarito.
                \item \textbf{CUIDADO ao preencher o gabarito} pois eles são \textbf{INDIVIDUAIS
                    e NOMEADOS} (cada estudante tem um gabarito próprio específico, com um
                    código de identificação único). \textbf{Questões rasuradas são anuladas}
                    pelo software de leitura do gabarito.
                \item Ao final da prova, \textbf{DEVOLVA A PROVA E O GABARITO} para o professor.
                \item Siga todas as normas de \textbf{Integridade Acadêmica} da disciplina.
                    Alunos que forem flagrados com qualquer espécie de ``cola'' ou trocando
                    informações com outros alunos terão suas avaliações recolhidas, as notas
                    zeradas, e a situação será encaminhada para a coordenação para a aplicação
                    das penalidades previstas pela universidade.
                \item Com 90 minutos de avaliação você terá tempo de até 2,min por questão,
                    em média.
                \item Boa avaliação!
            \end{itemize}
            \vspace{2cm}
        }}}
    \end{center}
    \newpage

    %%%%%%%%%%%%%%%%%%%%%%%%%%%%%%%%%%%%%%%%%%%%%%%%%%%%%%%%%%%%%%%%%%%%%%%%%%%%%%%%
    %%% Inicia o bloco de questões:
    \begin{questions}

    \newpage
    \question

(Adaptado de ENADE 2011) Considere o diagrama abaixo, que ilustra um banco de

dados que descreve os empregados que trabalham nos departamentos de uma empresa:

\begin{figure}[H]

  \centering

  \caption{Empregados e departamentos}

  \vspace{-0.2cm}

  \label{fig:empregados}

  %\fbox{

    \includegraphics[width=0.5\linewidth]{static/imgs/defaul.png}

  %}\\

  %\footnotesize{Fonte: xxxx.}

\end{figure}

\vspace{-0.4cm}

Na linguagem SQL, o comando mais simples para recuperar os códigos dos

departamentos que contém empregados que ganham menos do que 2000 e mais do que

5000 reais é:


\begin{choices}
\choice \begin{verbatim}SELECT cod_departamento

FROM empregados

WHERE salario <= 2000 AND salario >= 5000;

\end{verbatim}


\choice \begin{verbatim}SELECT cod_departametno

FROM empregados

WHERE salario <= 2000 OR salario >= 5000;

\end{verbatim}


\choice \begin{verbatim}SELECT cod_departamento

FROM empregados

WHERE salario < 2000 AND salario > 5000;

\end{verbatim}


\choice \begin{verbatim}SELECT cod_departamento

FROM empregados

WHERE salario < 2000 OR salario > 5000;

\end{verbatim}


\choice \begin{verbatim}SELECT cod_departamento

FROM empregados

WHERE salario BETWEEN 2000 AND 5000;

\end{verbatim}


\end{choices}

\question

Assinale quantas sequências de caracteres a seguir são reconhecidas pelo autômato finito abaixo. As quatro sequências de caracteres (separadas por vírgulas) são: \(0\), \(+567\), \(-89.5\), \(-3e3\).



\begin{figure} [H]

    \centering

    \includegraphics[width=0.5\linewidth]{static/imgs/defaul.png}

    \caption{Questão 25}

    \label{fig:enter-label}

\end{figure}



Qual das seguintes alternativas indica o número de sequências reconhecidas pelo autômato?
\begin{choices}
\choice 2
\choice 0
\choice 1
\choice 4
\choice 3
\end{choices}

\question

Em SQL, qual comando é utilizado para combinar linhas de duas ou mais tabelas?


\begin{choices}
\choice CONNECT
\choice MERGE
\choice JOIN
\choice LINK
\choice COMBINE
\end{choices}

\question

Assinale a alternativa INCORRETA:


\begin{choices}
\choice O controle de fluxo tem como objetivo garantir que nenhum dos parceiros de uma comunicação inunda o outro enviando pacotes mais rápido do que ele pode tratar.
\choice Nos serviços orientados a conexões há a necessidade de estabelecimento de uma conexão antes da transferência dos dados.
\choice Os serviços orientados a conexão podem ajudar no controle de congestionamento através da diminuição da taxa de transmissão durante um congestionamento em andamento.
\choice Os serviços orientados a conexões são sempre confiáveis garantindo a entrega ordenada e completa dos dados transmitidos.
\choice Serviços orientados a conexão podem ser implementados em subredes que funcionam no modo datagrama.
\end{choices}

\question

Um sistema centralizado é um concentrador de recursos; um sistema distribuído apresenta seus recursos dispersos. Entretanto, nem todo conjunto de recursos computacionais dispersos pode ser considerado um sistema distribuído. Considerando um conjunto de computadores, assinale a alternativa que melhor corresponde às características necessárias para considerá-lo um sistema distribuído:
\begin{choices}
\choice Suporte de rede e um relógio global
\choice Existência de sistema operacional idêntico e hardware padronizado em todos os computadores
\choice Suporte de rede e funções primitivas de comunicação 
\choice Existência de memória compartilhada e relógios locais sincronizados 
\choice Existência de memória secundária compartilhada e protocolos de sincronização de estado
\end{choices}

\question

Uma característica de uma arquitetura RISC é:
\begin{choices}
\choice Uma arquitetura de alto risco pois o mercado de hardware evolui muito rapidamente
\choice Possui um grande conjunto de instruções complexas
\choice A arquitetura é constituída de milhares de processadores
\choice O processador é formado por válvulas e transistores
\choice Possui um pequeno conjunto de instruções simples
\end{choices}

\question

A função abaixo, escrita na linguagem C, quando executada para \( n = 5 \), faz quantas chamadas recursivas (excluindo a primeira chamada da função)? 

\begin{verbatim} int fat (int n) 

{ 

if (n == 1) return n; 

else return (n*fat(n-1)); 

} 

\end{verbatim}
\begin{choices}
\choice 4
\choice 1
\choice 6
\choice 0
\choice 5
\end{choices}

\question

Em relação ao paradigma de programação cliente-servidor. Qual das afirmativas abaixo é FALSA?
\begin{choices}
\choice Um aplicativo servidor inicia ativamente o contato com clientes 
\choice Um aplicativo servidor aceita contato de clientes arbitrários, mas oferece um único serviço. 
\choice Um aplicativo cliente é um programa arbitrário que se torna temporariamente um cliente quando for necessário o acesso remoto a um serviço, mas também executa processamento local.
\choice Um aplicativo servidor é um programa de propósito especial dedicado a fornecer um serviço, mas pode tratar de múltiplos clientes remotos ao mesmo tempo. 
\choice Um aplicativo cliente pode acessar múltiplos serviços quando necessário.
\end{choices}

\question

Um agente SNMP é um aplicativo que é executado:
\begin{choices}
\choice em roteadores com filtragem de pacotes com o objetivo de proteger acesso a rede
\choice em "firewalls" com o objetivo de proteger acesso a rede
\choice a partir de um computador específico para monitorar a rede
\choice em computadores denominados de gerentes
\choice em um dispositivo de rede
\end{choices}

\question

O usuário que não é da área de tecnologia, costuma ser de alta gerência, que é

especialista no negócio e tem a responsabilidade central pelas políticas de

dados da empresa é o:
\begin{choices}
\choice Gerente de Tecnologia da Informação (ITM: IT Manager)
\choice Coordenador de Tecnologia da Informação (ITS: IT Supervisor)
\choice Encarregado da Proteção de Dados (DPO: Data Protection Officer)
\choice Administrador do Banco de Dados (DBA: Database Administrator)
\choice Administrador de Dados (DA: Data Administrator)
\end{choices}

\question

Qual das seguintes opções \textbf{não é} uma vantagem do uso de um banco de

dados sobre o sistema de arquivos convencional?


\begin{choices}
\choice Estruturação de dados
\choice Controle de acesso e visualização de dados
\choice Controle de redundância e inconsistência
\choice Independência de dados
\choice Eficiência
\end{choices}

\question

Considere um problema em que são dados 5 objetos com os seguintes pesos e valores:\\

\text{pesos: } $(W_1, W_2, W_3, W_4, W_5) = (6, 10, 9, 5, 12)$\\

\text{valores: } $(P_1, P_2, P_3, P_4, P_5) = (8, 5, 10, 15, 7)$\\

Além disso, é dada uma mochila que suporta até 30 unidades de peso, para transportar os objetos. O objetivo do problema é preencher a mochila de tal forma que o valor total dos objetos a serem transportados seja o maior possível, mas sem exceder o limite de peso suportado pela mochila. Assuma que é permitido colocar fração de um objeto na mochila. Qual das seguintes alternativas corresponde ao valor máximo obtido no preenchimento da mochila:
\begin{choices}
\choice 12.2
\choice 21.5
\choice 43.1
\choice 30.34
\choice 38.83
\end{choices}

\question

Todos os convidados presentes num jantar tomam chá ou café. Treze convidados bebem café, dez bebem chá e 4 bebem chá e café. Quantas pessoas tem nesse jantar.


\begin{choices}
\choice 23
\choice 19
\choice 10
\choice 15
\choice 27
\end{choices}

\question

Para cada \( n \in \mathbb{N} \), seja \( D_n = (0, 1/n) \), onde \( (0, 1/n) \) representa o intervalo aberto de extremos \( 0 \) e \( 1/n \). O conjunto diferença \( D_3 - D_{20} \) é igual a:


\begin{choices}
\choice \( D_{20} \cup D_3 \)
\choice \( D_3 \)
\choice \( [1/20, 1/3) \)
\choice \( (1/20, 1/3) \)
\choice \( D_{20} \)
\end{choices}

\question

Considere o seguinte trecho de programa:\\

1. \( i := 1; \)\\

2. while \( i \leq n \) do\\

   begin\\

3. \( sum := sum + a[i]; \)\\

4. \( i := i + 1; \)\\

   end;\\



Considere que:\\

- \( I \) representa a inicialização da variável \( i := 1 \) na linha 1;\\

- \( T \) representa o teste da linha 2;\\

- \( A \) representa os comandos da linha 3;\\

- \( P \) representa o incremento na linha 4.\\



Qual é a expressão regular que representa todas as sequências de passos possíveis de serem executados por este trecho de programa?
\begin{choices}
\choice \( I(TAP)^* \)
\choice \( I(TAP)^+ \)
\choice \( I(T(AP)^*)^* \)
\choice \( IT^+A^*P^* \)
\choice \( I(T(AP)^*)^+ \)
\end{choices}

\question

Seja \( f : \mathbb{R} \to \mathbb{R} \) derivável. Se existem \( a, b \in \mathbb{R} \) tal que \( f(a)f(b) < 0 \) e \( f(x) \neq 0 \) para todo \( x \in (a, b) \), podemos afirmar que no intervalo \( (a, b) \) a equação \( f(x) = 0 \) tem:
\begin{choices}
\choice somente raízes imaginárias
\choice uma única raiz real
\choice duas raízes reais
\choice uma raiz imaginária
\choice nenhuma raiz real
\end{choices}

\question

Qual dos protocolos abaixo pode ser caracterizado como protocolo de roteamento do tipo estado de enlace?
\begin{choices}
\choice IGMP
\choice ICMP
\choice RIP2
\choice OSPF
\choice BGP-4
\end{choices}

\question

Na criptografia com chave pública:


\begin{choices}
\choice O sigilo é obtido através da codificação com a chave pública do destinatário e decifragem com a chave privada do destinatário.
\choice Para assinar digitalmente uma mensagem codifica-se a mesma com a chave pública do remetente e esta é decifrada com a chave privada do destinatário.
\choice O sigilo é obtido através da codificação com a chave privada do remetente e decifragem com a chave pública do destinatário.
\choice Para assinar digitalmente uma mensagem codifica-se a mesma com a chave pública do destinatário e esta é decifrada com a chave privada do destinatário.
\choice O sigilo é obtido através da codificação com a chave privada do destinatário e decifragem com a chave pública do destinatário.
\end{choices}

\question

Na Álgebra Booleana:
\begin{choices}
\choice Os dígitos são hexadecimais de 0 a 15
\choice Os dígitos são alfanuméricos que podem ser representados por um byte
\choice Há dez valores motivados pelos dez dedos do ser humano
\choice Os dígitos são binários 0 e 1
\choice Os dígitos são octais, de 0 a 7
\end{choices}

\question

Pode-se afirmar que o gráfico da função \(y = 2 + \frac{1}{x - 1}\) é o gráfico da função \(y = \frac{1}{x}\):
\begin{choices}
\choice transladado uma unidade para a direita e duas unidades para baixo
\choice transladado uma unidade para a esquerda e duas unidades para cima
\choice transladado uma unidade para a direita e duas unidades para cima
\choice transladado uma unidade para a esquerda e duas unidades para baixo
\choice nenhuma das anteriores.
\end{choices}

\question

Histograma de uma imagem com \( K \) tons de cinza é:
\begin{choices}
\choice Contagem do número de vezes que cada um dos \( K \) tons de cinza ocorreu na imagem.
\choice Contagem do número de objetos encontrados na imagem.
\choice Contagem do número de tons de cinza que ocorreram na imagem.
\choice Contagem dos pixels da imagem.
\choice Nenhuma alternativa acima.
\end{choices}

\question

Marque a alternativa onde todos os conceitos estão corretos.


\begin{choices}
\choice Em um diagrama de fluxo de dados uma entidade externa representa um produtor ou um consumidor de informação e está fora dos limites do sistema modelado; cada processo deve ser refinado, para explicitar um maior detalhamento; um DFD pode conter vários níveis de detalhamento; um processo é um transformador de informação e também está fora do sistema; o nível 0 de um DFD representa o sistema como um todo e indica os principais usuários e as funções do sistema.
\choice Em um diagrama de fluxo de dados uma entidade externa representa uma fonte ou destino das informações processadas pelo sistema e está fora dos limites do sistema modelado; cada processo pode ser refinado, para explicitar um maior detalhamento; um DFD pode conter vários níveis de detalhamento; um processo é um transformador de informação e também está fora do sistema; o nível 0 de um DFD representa o sistema como um todo e indica as principais fontes e destinos das informações.
\choice Em um diagrama de fluxo de dados, uma entidade externa representa um produtor ou um consumidor de informação e está fora dos limites do sistema modelado; cada processo pode ser refinado, para explicitar um maior detalhamento; um DFD contém dois níveis de detalhamento; um processo é um transformador de informação e também está fora do sistema; o nível 0 de um DFD representa o sistema como um todo e indica os principais usuários e as funções do sistema.
\choice Nenhuma das alternativas anteriores.
\choice Em um diagrama de fluxo de dados uma entidade externa representa uma fonte ou destino das informações processadas pelo sistema e está fora dos limites do sistema modelado; cada processo pode ser refinado, para explicitar um maior detalhamento; um DFD pode conter vários níveis de detalhamento; um processo é um transformador de informação; o nível 0 de um DFD representa o sistema como um todo e indica as principais fontes e destinos das informações, usualmente referenciado por Diagrama de Contexto.
\end{choices}

\question

Considere as seguintes tabelas em uma base de dados relacional (chaves primárias sublinhadas):\\

Departamento (\underline{CodDepto}, NomeDepto)\\

Empregado (\underline{CodEmp}, NomeEmp, CodDepto)\\

Considere as seguintes restrições de integridade sobre esta base de dados relacional:\\

- Empregado.CodDepto é sempre diferente de NULL\\

- Empregado.CodDepto é chave estrangeira da tabela Departamento com cláusulas ON DELETE RESTRICT e ON UPDATE RESTRICT\\

Qual das seguintes validações não é especificada por estas restrições de integridade:
\begin{choices}
\choice Sempre que uma nova linha for inserida em Departamento, deve ser garantido que o valor de Departamento.CodDepto aparece na coluna Empregado.CodDepto.
\choice Sempre que o valor de Empregado.CodDepto for alterado, deve ser garantido que o novo valor de Empregado.CodDepto aparece em Departamento.CodDepto.
\choice Sempre que uma nova linha for inserida em Empregado, deve ser garantido que o valor de Empregado.CodDepto aparece na coluna Departamento.CodDepto.
\choice Sempre que o valor de Departamento.CodDepto for alterado, deve ser garantido que não há uma linha com o antigo valor de Departamento.CodDepto na coluna Empregado.CodDepto.
\choice Sempre que uma linha for excluída de Departamento, deve ser garantido que o valor de Departamento.CodDepto não aparece na coluna Empregado.CodDepto.
\end{choices}

\question

Três atletas \( A \), \( B \) e \( C \) competiram, ao pares, numa corrida de \( d \) metros. Considerando que cada atleta teve o mesmo desempenho (ou seja, a mesma velocidade) ao competir com adversários distintos, e sabendo-se que



\begin{itemize}

    \item \( A \) venceu \( B \) chegando 20 metros à frente

    \item \( B \) venceu \( C \) chegando 10 metros à frente

    \item \( A \) venceu \( C \) chegando 28 metros à frente,

\end{itemize}



podemos afirmar que a corrida tem
\begin{choices}
\choice 150 metros
\choice 200 metros
\choice 50 metros
\choice 110 metros
\choice 100 metros
\end{choices}

\question

Ao analisar um novo sistema de gerenciamento de bancos de dados que surgiu no

mercado, você percebe que a entrada às operações do modelo de dados são tabelas

mas as saídas nunca são tabelas, são sempre outras estruturas. O vendedor desse

SGBD afirma que ele é um SGBD relacional. Pode-se afirmar que:


\begin{choices}
\choice Esse SGBD não é relacional pois no modelo relacional a única estrutura percebida pelos usuários são tabelas.
\choice Esse SGBD não é relacional pois, para ser verdadeiramente relacional, as operações do modelo de dados não podem trabalhar recebendo tabelas como entrada.
\choice Esse SGBD é relacional pois consegue trabalhar com tabelas para a entrada das operações do modelo de dados. A saída não é importante.
\choice Nenhuma das alternativas anteriores.
\choice Esse SGBD é relacional pois o vendedor afirmou que é.
\end{choices}

\question

Qual das seguintes afirmações sobre crescimento assintótico de funções não é verdadeira:
\begin{choices}
\choice \( 2n^2 + 3n + 1 = O(n^2) \)
\choice \( \log n^2 = O(\log n) \)
\choice Se \( f(n) = O(g(n)) \) e \( g(n) = O(h(n)) \) então \( f(n) = O(h(n)) \)
\choice \( 2^{n+1} = O(2^n) \)
\choice Se \( f(n) = O(g(n)) \) então \( g(n) = O(f(n)) \)
\end{choices}

\question

Se \( \lim_{n \to \infty} \left( \frac{1}{n^2} + \frac{2}{n^2} + \cdots + \frac{n-1}{n^2} \right) = L \) então


\begin{choices}
\choice \( L = 2 \)
\choice \( L = \infty \)
\choice \( L = 0 \)
\choice \( L = 1 \)
\choice \( L = 1/2 \)
\end{choices}

\question

Considere a seguinte tabela para uma base de dados relacional:\\

Empregado (CodEmp, NomeEmp, CodDepto)\\

Considere que esta tabela tem um índice na forma de uma árvore B sobre as colunas (CodEmp, CodDepto), nesta ordem.

Quanto a este índice, considere as seguintes afirmativas:

\begin{enumerate}

    \item Este índice pode ser usado pelo SGBD relacional para acelerar uma consulta na qual são fornecidos os valores de CodEmp e CodDepto.

    \item Este índice pode ser usado pelo SGBD relacional para acelerar uma consulta na qual é fornecido um valor de CodEmp.

    \item Este índice não é adequado para ser usado pelo SGBD relacional para acelerar uma consulta na qual é fornecido um valor de CodDepto.

    \item O algoritmo que faz inserções e remoções de entradas do índice tem por objetivo garantir que o índice fique organizado de tal forma que o acesso a cada nodo da árvore implique em número de acessos semelhantes.

    \item O índice por árvore-B não é adequado para tabelas que sofrem grande número de inclusões e exclusões, pois exige reorganizações frequentes.

\end{enumerate}





Quanto a estas afirmativas pode se dizer que:
\begin{choices}
\choice Nenhuma das afirmativas está correta
\choice Apenas as afirmativas 1), 2) e 4) estão corretas
\choice Apenas as afirmativas 1), 2) e 5) estão corretas
\choice Todas afirmativas estão corretas
\choice Apenas as afirmativas 1), 2), 3) e 4) estão corretas
\end{choices}

\question

A média aritmética de uma lista de 50 números é 50. Se dois desses números, 51 e 97, forem suprimidos dessa lista a média dos restantes será


\begin{choices}
\choice 49
\choice 40
\choice 50
\choice 47
\choice 51
\end{choices}

\question

Dado um vetor \( u \in \mathbb{R}^2 \), \( u = (-3, 4) \), vamos denotar por \( v \) o vetor de \( \mathbb{R}^2 \) que tem tamanho 1 e é ortogonal a \( u \). Então \( v \) pode ser dado por


\begin{choices}
\choice \( \left( -\frac{4}{5}, -\frac{3}{5} \right) \)
\choice \( \left( -\frac{4}{5}, \frac{1}{5} \right) \)
\choice \( \left( \frac{3}{5}, \frac{4}{5} \right) \)
\choice \( \left( -\frac{4}{5}, \frac{3}{5} \right) \)
\choice \( \left( -\frac{4}{5}, \frac{2}{5} \right) \)
\end{choices}

\question

Os Sistemas de Bancos de Dados (SBD) são essenciais hoje em dia, tanto para

pessoas físicas quanto para pessoas jurídicas. Assinale a alternativa que não

representa adequadamente a importância do uso de um SBD:
\begin{choices}
\choice Permitem consultas históricas, de modo otimizado, para que alguma informação atual possa ser comparada com dados anteriores.
\choice Permitem que as empresas tenham mecanismos de obedecer à Lei Geral de Proteção de Dados Pessoais (LGPD), impedindo processos de anonimização das informações.
\choice Permitem às empresas coletar, armazenar, agregar, manipular e gerenciar dados, tomando assim melhores decisões de negócios e fornecendo diversos serviços úteis aos seus clientes.
\choice Permitem que as empresas criem sites na Internet de modo que os próprios usuários consigam consultar e atualizar alguns dados.
\choice ermitem que as empresas criem aplicativos para telefones celulares, oferecendo serviços aos seus usuários.
\end{choices}

\question

Quando dizemos que um banco de dados é em \ingles{cloud} e centralizado, estamos

nos referindo, respectivamente, à localização dos dados e da

infraestrutura. Verdadeiro ou falso?
\begin{choices}
\choice Não é possível determinar com os dados apresentados.
\choice Falso
\choice Verdadeiro
\end{choices}

\question

Considere as afirmações abaixo sobre endereços IP:

\begin{itemize}

    \item[I.] Uma rede IP classe C fornece até 256 endereços válidos para serem atribuídos a equipe.

    \item[II.] A quantidade máxima de bits que pode ser utilizada para se definir sub-redes em uma rede IP classe C é seis (6).

    \item[III.] A máscara padrão para uma rede classe B é 255.255.255.0.

\end{itemize}

Qual das seguintes alternativas representa as afirmações corretas?
\begin{choices}
\choice Somente I e II.
\choice Somente III.
\choice Somente II.
\choice Somente I.
\choice Somente II e III.
\end{choices}

\question

São vantagens da utilização de threads no espaço do usuário, exceto:
\begin{choices}
\choice Maior portabilidade da aplicação.
\choice A criação e o gerenciamento das threads são mais eficientes.
\choice O sistema operacional escalona a thread.
\choice O escalonamento pode ser específico para a aplicação.
\choice Nenhuma modificação é necessária no kernel.
\end{choices}

\question

Um compilador detecta:
\begin{choices}
\choice erros nos resultados gerados pelo programa
\choice todos os erros citados acima
\choice erros que podem ocorrer durante a execução do programa
\choice erros aritméticos
\choice erros de sintaxe do programa 
\end{choices}

\question

Dentre as definições a seguir, ligadas ao conceito de normalização do modelo relacional, qual delas é \textbf{INCORRETA}?


\begin{choices}
\choice A normalização é um processo passo a passo reversível de substituição de uma dada coleção de relações por sucessivas coleções de relações as quais possuem uma estrutura progressivamente mais simples e mais regular.
\choice A primeira forma normal estabelece que os atributos da relação contêm apenas valores atômicos.
\choice O objetivo da normalização é eliminar várias anomalias (ou aspectos indesejáveis) de uma relação.
\choice As formas normais se baseiam em certas estruturas de dependências.
\choice As relações que obedecem à primeira forma normal não apresentam anomalias.
\end{choices}

\question

Considerando a rede de Petri abaixo, quais das alternativas são verdadeiras?



\begin{itemize}

    \item[I)] O lugar \( A \) está habilitado a disparar.

    \item[II)] Apenas a transição \( T1 \) está habilitada a disparar.

    \item[III)] A sequência de transições \( (T1, T2, T3, T2) \) pode ser disparada, nessa ordem.

    \item[IV)] A transição \( T4 \) nunca poderá ser disparada.

\end{itemize}



\begin{figure}[H]

    \centering

    \includegraphics[width=0.5\linewidth]{static/imgs/defaul.png}

    \caption{Questão 49}

    \label{fig:enter-label}

\end{figure}
\begin{choices}
\choice Apenas as alternativas I e III.
\choice Apenas as alternativas II e III.
\choice Apenas as alternativas II, III e IV.
\choice Todas as alternativas.
\choice Apenas as alternativas III, IV.
\end{choices}

\question

Que cláusula da SQL é fortemente relacionada com a lógica matemática?


\begin{choices}
\choice WHERE
\choice INNER JOIN
\choice ORDER BY
\choice SELECT
\choice FROM
\end{choices}

\question

Dentre as definições a seguir, ligadas ao conceito de visões do modelo relacional, qual delas é \textbf{INCORRETA}?


\begin{choices}
\choice Programas aplicativos do banco de dados podem ser executados sobre visões de relações da base de dados.
\choice Uma visão relacional é uma relação virtual, derivada de relações base a partir da especificação de operações da álgebra relacional.
\choice O gerenciamento de visões envolve a conversão da consulta do usuário sobre as visões para a consulta sobre as relações base.
\choice Uma visão relacional é uma relação virtual que nunca é materializada.
\choice Uma visão é útil por representar uma percepção particular do banco de dados, compartilhado por muitos aplicativos.
\end{choices}

\question

Quais das seguintes afirmações são verdadeiras? As Métricas de software servem para:



\begin{itemize}

    \item[I)] indicar a qualidade do produto e avaliar a produtividade.

    \item[II)] auxiliar na melhoria do processo.

    \item[III)] formar uma base para as estimativas e justificar a aquisição de ferramentas.

    \item[IV)] determinar se a utilização de um método traz benefícios ou não.

\end{itemize}
\begin{choices}
\choice Apenas as alternativas II e III.
\choice Apenas as alternativas I, II e IV.
\choice Nenhuma delas.
\choice Todas as alternativas.
\choice Apenas as alternativas I, IV.
\end{choices}

\question

A velocidade de um ponto em movimento é dada pela equação\\

$v(t) = t e^{-0.01t} \, \text{m/s}$\\

O espaço percorrido desde o instante que o ponto começou a se mover até a sua parada total é


\begin{choices}
\choice \( 10^2 e^{-1} \, m \)
\choice \( 10^4 \, m \)
\choice \( 10^3 e^{-0.01} \, m \)
\choice \( (e^{-100} - 1) m \)
\choice \( 10^2 \, m \)
\end{choices}

\question

Quais dos algoritmos de ordenação abaixo possuem tempo no pior caso e tempo médio de execução proporcional a \( O(n \log n) \).
\begin{choices}
\choice Quicksort e merge sort
\choice Merge sort e bubble sort
\choice Merge sort e heap sort
\choice Heap sort e selection sort
\choice Bubble sort e Quick sort
\end{choices}

\question

Um banco de dados (BD) é uma coleção logicamente coerente de dados relacionados

e persistentes. Em relação às características de um banco de dados, assinale a

alternativa verdadeira:
\begin{choices}
\choice Geralmente é projetado e usado para finalidades gerais, sendo extremamente genéricos e adaptáveis para qualquer situação.
\choice Não precisa de dados atuais e verdadeiros para ser útil (para representar corretamente a realidade atual).
\choice Não precisam ser obrigatoriamente computadorizados, existem bancos de dados manuais e/ou em papel.
\choice Armazena dois grandes ``tipos`` de dados: os dados do usuário final e os metadados, sendo que os metadados ficam armazenados nas mesmas tabelas que os dados dos usuários e isso nos permite saber a estrutura dos bancos de dados.
\choice Representa um aspecto (ou parte) do mundo real, parte essa chamada de ``catálogo' ou ``dicionário de dados``.
\end{choices}

\question

O pipeline de visualização de objetos tridimensionais reúne um conjunto de transformações e processos aplicados a primitivas geométricas. Sobre essas transformações e processos pode-se dizer que:

\begin{enumerate}[label=\Roman*.]

    \item Os objetos devem corresponder a sólidos.

    \item As coordenadas dos vértices sofrem transformação de acordo com a posição e orientação do observador.

    \item Um volume de visualização correspondente a um paralelepípedo é determinado pela adoção de projeção perspectiva.

    \item A fase final do pipeline corresponde à rasterização dos polígonos.

\end{enumerate}

Selecione a alternativa correta:
\begin{choices}
\choice Apenas a afirmativa IV está verdadeira.
\choice Apenas as afirmativas I e III são falsas.
\choice Apenas a afirmativa IV é falsa.
\choice Todas as afirmativas são verdadeiras.
\choice As afirmativas II e III são falsas.
\end{choices}

\question

Quanto ao TCP, é INCORRETO afirmar:
\begin{choices}
\choice Utiliza portas para permitir a comunicação entre processos localizados em dispositivos diferentes.
\choice Possui um campo de \textit{checksum} que valida as informações de seu cabeçalho, mas não valida as informações de \textit{payload} (campo de dados).
\choice É um protocolo do nível de transporte.
\choice É um protocolo orientado a conexão.
\choice Usa janelas deslizantes para implementar o controle de fluxo e erro.
\end{choices}

\question

O número de \textit{strings binárias} de comprimento 7 e contendo um par de zeros consecutivos é


\begin{choices}
\choice 94
\choice 90
\choice 92
\choice 91
\choice 95
\end{choices}

\question

Numa prova de múltipla escolha com 10 questões e 4 alternativas, qual a chance (probabilidade) de um aluno apenas "chutando as respostas" conseguir "gabaritar" a prova (acertar todas as questões).


\begin{choices}
\choice \( \frac{1}{4^{15}} \)
\choice \( \frac{1}{10^{8}} \)
\choice \( \frac{1}{10^{4}} \)
\choice \( \frac{1}{4^{20}} \)
\choice \( \frac{1}{2^{20}} \)
\end{choices}

\question

Analise o diagrama relacional do banco de dados exibido na figura abaixo:

\begin{figure}[H]

  \centering

  \caption{Banco de dados de peças e fornecedores}

  \label{fig:pecas}

  %\fbox{

    \includegraphics[width=0.5\linewidth]{static/imgs/defaul.png}

  %}\\

  %\footnotesize{Fonte: xxxx.}

\end{figure}

Podemos dizer que esse banco de dados ilustra:
\begin{choices}
\choice Um relacionamento com cardinalidade N:N entre as tabelas ``pecas' e ``fornecedores', realizado através de relacionamentos identificados através da tabela ``fornecimentos'. 
\choice Um relacionamento com cardinalidade 1:N entre as tabelas ``pecas' e ``fornecedores', indicando que uma peça pode ter sido fornecida por diversos fornecedores. 
\choice Um relacionamento com identificação N:N entre as tabelas ``pecas' e ``fornecedores', realizado através do relacionamento com cardinalidade identificada através da tabela ``fornecimentos'. 
\choice Um relacionamento com cardinalidade N:N entre as tabelas ``pecas' e ``fornecimentos', realizado através de relacionamentos não identificados com a tabela ``fornecedores'. 
\choice Um relacionamento com cardinalidade N:N entre as tabelas ``pecas' e ``fornecedores', realizado através de relacionamentos não identificados através da tabela ``fornecimentos'. 
\end{choices}

\question

Um sistema de gerenciamento de bancos de dados (SGBD) é um software servidor

que nos permite definir, construir, manipular, integrar e compartilhar bancos

de dados entre diversos usuários e aplicações. São exemplos de

SGBD, \textbf{exceto}:
\begin{choices}
\choice PostgreSQL
\choice Oracle SQL Developer
\choice MongoDB
\choice Oracle Database
\choice SQL Server
\end{choices}

\question

Seja a árvore binária abaixo a representação de um espaço de estados para um problema \( p \), em que o estado inicial é \( a \), e \( i \) e \( f \) são estados finais.



\begin{figure}[H]

    \centering

    \includegraphics[width=0.5\linewidth]{static/imgs/defaul.png}

    \caption{Questão 59}

    \label{fig:enter-label}

\end{figure}



Um algoritmo de busca em largura-primeiro forneceria a seguinte sequência de estados como primeira alternativa a um caminho-solução para o problema \( p \):
\begin{choices}
\choice \( a \ b \ c \ d \ e \ f \)
\choice \( a \ b \ e \ i \)
\choice \( a \ b \ d \ h \ e \ i \)
\choice \( a \ b \ d \ e \ f \)
\choice \( a \ c \ f \)
\end{choices}



    %%%%%%%%%%%%%%%%%%%%%%%%%%%%%%%%%%%%%%%%%%%%%%%%%%%%%%%%%%%%%%%%%%%%%%%%%%%%%%%%
    %%% Finaliza o bloco de questões:
    \end{questions}
    \end{document}
    